
\AddSymbolsA
\pagecolor{OliveGreen!30!gray}
\color{white}
\quad\quad \lettrine[,lraise=0.2,lhang=0.5]{\textbf{街}}{道}上传来伊藤佐一\index{yi@伊藤佐一}的脚步声。那脚步声，是细碎的——却足以吓退草丛中驻足的一只猫\index{mao@猫}；猫转过头来。

伊藤佐一看着那只橘黄猫的纺锤形瞳孔——他猛地瞧见里面站着一个人。

\twkkai{是我自己吗？}他探下身，伸手向那只猫挥了一挥。

猫不作声，却也不动。

伊藤佐一扭头看向四周；行人们都自顾自地走着。他们的无聊反衬着伊藤佐一对这只猫的好奇。他伸手从自己拿着的纸袋中掏出一块面包，撕成一块块，开始逗起了猫——他并不知道猫对面包感不感兴趣。

猫出乎意料地从草丛走出，前脚掌撑在街道的粗糙水泥地上，低头嗅了嗅那些面包屑，却并不吃，抬头望向了伊藤佐一。

伊藤佐一这才留意到，猫的颈部挂着一块精巧的金质奖章。

他缓慢地伸出手，却见那只猫也开始伸出猫爪，轻轻搭在了他的掌心。伊藤佐一却也不管，用左手直接托起猫，右手捏紧了那块奖章。

猫开始嚎叫起来。

而他看见了什么：
\begin{gather*}
L(s,\chi)=\sum_{n=1}^{\infty}\frac{\chi(n)}{n^s}=\prod_{p}\left(1-\frac{\chi(p)}{p^s}\right)^{-1},\  \Re(s)>1.\\
1837.
\end{gather*}\index{lhanshu@$L(s,\chi)$（狄利克雷$L$-函数）!定义式: $L(s,\chi) = \sum_{n=1}^{\infty} \chi(n)/n^s\ (\Re(s)>1)$}\index{lhanshu@$L(s,\chi)$（狄利克雷$L$-函数）!欧拉乘积分解：$L(s,\chi)=\prod_{p}\left(1-\chi(p)p^{-s}\right)^{-1}\ (\Re(s)>1)$}
\ \ \ \ \ \twkkai{这是……？}

\twkkai{喵！}

猫挣脱了伊藤佐一的手，开始狂奔，又躲进了草丛里。伊藤佐一仿佛还能听到猫的喘息声。

那块奖章意外地掉落，潜伏在面包屑之中——自然躲不过伊藤佐一的目光。他拾起它，犹豫着——是戴回猫脖子上，还是据为己有。

长时间的精神挣扎让他感到疲惫；他决定再看那只猫一眼——它早已不见踪影。

\twkkai{看来只能自私点了；}伊藤佐一看着奖章，又沉思着。

\twkkai{1837……明天问问教授\index{jiao@教授}吧；}伊藤佐一把奖章揣进了口袋。

\begin{center}
——————
\end{center}

第二天，伊藤佐一穿过嬉闹的人群，走进了课室。

那课室，曾经举办过风格各异的讲座，每一个不同的教授都曾有机会在这里大放异彩。一排排弧形的座位环绕着讲台，层层叠叠地分布在课室的中央与两端。

伊藤佐一坐下了。

他依旧惦记着那块奖章。他把它小心翼翼地拿出来，闻着——尽管奖章没有生锈，他的握着奖章的掌心依旧弥漫着一股金属味。

他盯着那几个符号——它们仿佛神谕中的咒符一般，与他的眼睛连成钢丝般柔韧的线，让他的目光迟迟不能转移：
\[\mathlarger{\mathlarger{\mathlarger{\Sigma}}}, \quad  \mathlarger{\mathlarger{\mathlarger{\Pi}}}, \quad \mathlarger{\mathlarger{\mathlarger{\chi}}}, \quad \mathlarger{\mathlarger{\mathlarger{\Re}}}\]

\twkkai{教授好像提及过这些符号——看来并不是新概念；}他沉静地想着，额头抵着桌面边沿。

这时，一只猫从课桌底下钻了进来。

正是那只猫。

它灵活地扭动着身躯，在伊藤佐一的裤腿间焦虑地跑动着，画出“\(\infty\)”形状的运动轨迹。

伊藤佐一却很迟才注意到这只猫。他突然嚷叫起来。

教授停止了讲课。伊藤佐一挺直了腰板，仿佛想把自己伪装进认真听课的人群中。

可是教授已经走到了伊藤佐一的身旁。伊藤佐一看向地面；猫已不见踪影。他尴尬地看向教授。旁边的同学开始笑了起来。

\twkkai{你在叫什么？你最近的状态很不对劲喔……你刚才在干什么？}

教授的声音虽然低沉柔和，但仍显得咄咄逼人。

\twkkai{我在想一些问题……和您讲的有关……}

\twkkai{嗯？我刚才在讲什么——你知道吗？}

伊藤佐一突然不说话了，又看向了黑板；可是他看不清。他开始紧张了，手心开始出汗了。他是那么害怕——害怕奖章掉落——以至于他开始把手伸进裤袋里。

\twkkai{你手里拿的是什么……我倒是挺好奇……；}教授轻轻地跺着脚。

众人的目光压迫着伊藤佐一。

他只好犹疑地拿出奖章，递给了教授。

教授一开始不以为意；可是后来，他的眼神愈发显露出光彩。

过了一会儿，他便转变了态度。

\twkkai{同学们！我们要感谢这位——你叫什么名字？}

\twkkai{伊藤佐一……}

\twkkai{很好，感谢这位同学；今天，我们将要拓展教材，讲讲我最近的研究方向。那么，作为第一课，讲课的主题就先定为……}

教授急急地走下阶梯，重新来到黑板前，写下几个伊藤佐一终于看得见的大字——“\(L\)-函数及其素因子分解”。

\twkkai{同学们；}教授举起粉笔，指着那些神秘的符号。
\begin{align*}
\sum_{n=1}^{\infty}\frac{\chi(n)}{n^s}&\longleftrightarrow\prod_{p}\left(1-\frac{\chi(p)}{p^s}\right)^{-1}\\
\sum_{n}&\longleftrightarrow\prod_{p}
\end{align*}

\twkkai{此时我已经写下了两个更简洁的符号：\(\sum_{n}\)和\(\prod_{p}\)，大家会有何联想？}

\twkkai{左边跑遍所有正整数\(n\)，右边跑遍所有素数\(p\)……；}伊藤佐一的直觉告诉他一个显然的结论。

\twkkai{没错，就是……}
\begin{align*}
\boxed{
n=\prod_pp^{e_p},\quad \text{\(e_p\geq0\), \twkkai{且只有有限个}\(e_p>0\).}}
\end{align*}

\twkkai{\(n\)的唯一素因子分解！}\index{weiyisuyinzifenjie@唯一素因子分解}伊藤佐一脑海中的声音和教授的高呼声重合，让他的脑袋嗡嗡地响着。

\twkkai{想象一下，\(\sum_n\)和\(\prod_p\)分别配备上\(\chi(n)\)和\(\chi(p)\)——就像两款不同的吸尘器，配备上装有不同滤网的真空袋……}\index{te@特征标$\chi$}

教授抬了抬眼镜，看着那些分神的学生，敲了敲黑板。

\twkkai{盯紧点！你们学到后面可就听不见如此生动的诠释了！你们要懂得一个道理——越抽象的真理越讨厌简陋的比喻；因此，不是我不给你们机会——是创造真理的上帝不给你们机会啊！}

同学们又诧异地看回教授，仿佛觉得他开始成为了一个自言自语的疯子，开始叽喳地讨论起来。

只有伊藤佐一仍在高速思考着。\twkkai{可是，\(\chi(n)\)和\(\chi(p)\)究竟有什么魔力呢？}

\twkkai{注意，重点来了；}教授仿佛在单独给伊藤佐一授课，每一句话都回应着他脑海里的声音。

教授的粉笔在黑板上嘶叫着，把自己的一部分化为神秘的算符。
\begin{align*}
\chi(mn)=\chi(m)\chi(n),\ \ 
\chi(p^k)=\chi(p)^k.
\end{align*}\index{te@特征标$\chi$!积性 $\chi(mn)=\chi(m)\chi(n)$}\index{te@特征标$\chi$!完全积性 $\chi(p^k)=\chi(p)^k$}
\ \quad \twkkai{这就是我所说的滤网——积性和完全积性性质；它们将正整数切割、分类……；}教授耐心十足地说着。

可是教授的讲解始终解决不了伊藤佐一心中弥漫着的疑惑。

\twkkai{为什么偏偏是那只猫呢？奖章……又意味着什么？还有最重要的，为什么奖章上刻的东西，在这个世界，竟然是可理解的？或者更准确地——为什么这些可理解之物，会被刻在不可理解的奖章上？}伊藤佐一左手托着下巴，手肘撑在膝盖上。而他的眼睛，目不转睛地盯着矩形的黑板：
\begin{gather*}
\left(1 - \frac{\chi(p)}{p^{s}}\right)^{-1}
= \sum_{k = 0}^{\infty} \left(\frac{\chi(p)}{p^{s}}\right)^{k}= \sum_{k = 0}^{\infty} \frac{\chi(p^{k})}{p^{ks}}.\\
(\text{\twkkai{这里用到了完全积性性质:} }\chi(p^k)=\chi(p)^k.)
\end{gather*}

\twkkai{等等，这里的动机究竟是什么？为什么突然引入几何级数？难道是……？}伊藤佐一很想举起他的手，但他不敢——尽管他坐在前排。

可是教授并没有停下来解释，继续写着；课室里飘荡着一阵死寂，只有粉笔和黑板的撞击声仍存。
\begin{gather*}
\text{\twkkai{由于}}\quad\sum_{k=0}^{\infty} \left| \frac{\chi(p^k)}{p^{ks}} \right|
    \leq \sum_{k=0}^{\infty} \frac{1}{p^{k \Re(s)}}
    = \frac{1}{1 - p^{-\Re(s)}}\quad\text{\twkkai{且}}\quad\prod_{p} \frac{1}{1 - p^{-\Re(s)}} = \zeta(\Re(s)) < \infty,\\
    (\text{\twkkai{这里用到了绝对收敛性.}})\\
\text{\twkkai{所以乘积}}\quad\prod_{p} \left( \sum_{k=0}^{\infty} \frac{\chi(p^{k})}{p^{ks}} \right)\quad\text{\twkkai{可以有效展开}}.
\end{gather*}
\index{juedui@绝对收敛性}



\twkkai{好了，同学们——这一切的一切，都是为了使用我们的终极武器——\(n\)的唯一素因子分解；它是让两款“吸尘器”都能正常运作的核心部件；}教授耐人寻味地说着，似乎暗中回答着伊藤佐一脑海中的私语。

\twkkai{那……怎样才能展开\(\prod_{p} \left( \sum_{k=0}^{\infty} \frac{\chi(p^{k})}{p^{ks}} \right)\)这个乘积呢？这可涉及到有限和无限的交融啊……}

伊藤佐一转着笔，又突然想到一件事。

\twkkai{那块奖章，教授不会忘了还给我吧？}

他决定，必须抢在下课前收拾好所有东西，拦住教授。

\twkkai{那么，在这里，我们要耍一个聪明的小花招：假设我们只考虑前 \(r\) 个素数 \(p_1, p_2, \dots, p_r\)，那么乘积就是有限个级数的乘积：}

教授在黑板上飞快地写着：
\[
\prod_{i=1}^{r} \left( \sum_{k=0}^{\infty} \frac{\chi(p_i^k)}{p_i^{k s}} \right).
\]

\twkkai{想象一下，把乘积展开的所有大括号里，每一个括号中的级数都有无限项；由于之前在黑板上我们已经证明过级数绝对收敛，那么，我们可以尝试像多项式乘法一样把它们逐项展开：  
从第一个级数中选一项，长这样： \(\frac{\chi(p_1^{e_1})}{p_1^{e_1 s}}\)，  
从第二个级数中选一项， \(\frac{\chi(p_2^{e_2})}{p_2^{e_2 s}}\)，……，  
再从第 \(r\) 个级数中选一项， \(\frac{\chi(p_r^{e_r})}{p_r^{e_r s}}\)，  
试着将它们相乘，我们似乎会得到一堆乱七八糟难以处理的式子：}
\[
\frac{\chi(p_1^{e_1})\chi(p_2^{e_2})\cdots\chi(p_r^{e_r})}{p_1^{e_1 s} p_2^{e_2 s}\cdots p_r^{e_r s}}.
\]

教授用粉笔狠狠地戳了戳黑板。

\twkkai{这里我们就要用到刚才的“滤网”和“终极武器”了，这便是化腐朽为神奇的重要一步：
由于 \(\chi\) 是完全积性的，分子的乘积等于 \(\chi(p_1^{e_1}p_2^{e_2}\cdots p_r^{e_r})\)。  
分母等于 \((p_1^{e_1}p_2^{e_2}\cdots p_r^{e_r})^{s}\)。
所以这一项就是 \(\frac{\chi(N)}{N^s}\)，其中 \(N = p_1^{e_1}p_2^{e_2}\cdots p_r^{e_r}\)。}

\twkkai{把每一个非负整数指数 \(e_1, e_2, \dots, e_r\) 分别对应到前 \(r\) 个素数，那么按照乘法 \(N = p_1^{e_1} p_2^{e_2} \cdots p_r^{e_r}\) 得到的 \(N\)，恰好跑遍所有只由前 \(r\) 个素数相乘得到的正整数（即素因子都在前 \(r\) 个素数中），并且不同的指数组 \((e_1, e_2, \dots, e_r)\) 会得到不同的 \(N\)，反过来每个这样的 \(N\) 也唯一对应一组指数。}

\twkkai{因此，我们此时会得到一个有趣的结果……}

教授粉笔的沙沙声和讲台下学生书写笔记的嘶嘶声互相映衬着。

而伊藤佐一远远地便望见了，那只猫，正坐在讲台上，用尾巴指着黑板的白色字迹：
\[
\prod_{i=1}^{r} \left( \sum_{k=0}^{\infty} \frac{\chi(p_i^k)}{p_i^{k s}} \right)
= \sum_{\substack{N \ge 1 \\ \text{\twkkai{所有素因子}} \le p_r}} \frac{\chi(N)}{N^s}.
\]

\twkkai{当 \(r \to \infty\)，左边就是原来的无穷乘积，右边则一定会跑遍所有正整数，因为每个正整数只有有限个素因子。于是形式上……}
\[
\prod_{p} \left( \sum_{k=0}^{\infty} \frac{\chi(p^k)}{p^{k s}} \right)
= \sum_{n=1}^{\infty} \frac{\chi(n)}{n^s}.
\]

教授没有意识到，学生们已经因为那只猫开始嬉笑起哄了。最前排的几个女同学正挤着嘴，试图吸引猫的注意力。

没曾想，教授一转过身来，猫便扑在他的身上，衔走了他手里的奖章，跑出了门外。

这一切，伊藤佐一自然看在眼里。\twkkai{我的奖章！}他的内心叫喊着，但他自然是无能为力的。

……

下课了，同学们纷纷离开课室——没有一个人在讨论刚才的数学问题。而教授则趁伊藤佐一经过时，拉住了他。

\twkkai{我一直觉得，你是我们班里最勤奋的学生……我希望经过今天这件事，你以后能多点来单独找我——我的办公室就在最顶层；}教授意味深长地对伊藤佐一说着。

伊藤佐一点了点头。

\twkkai{教授，你的研究方向是什么？}伊藤佐一没能忍住他的好奇。

\twkkai{类域论；到了某个阶段你就会明白，你今天所学的这一切，和这座大山联系是有多深的了；}\index{leiyulun@类域论（Class field theory）}教授拍了拍伊藤佐一的右肩。

\twkkai{走吧——我也要回去工作了；}教授整理好讲义，关了课室的灯，和伊藤佐一一起走出了门口。他们在门口互相告别了。

……

伊藤佐一独自走在马路上，仍为了他的奖章感到心心念念。他想不懂今天的数学课，想不懂类域论是什么，更想不懂它们之间的联系。

一辆车突然径直冲了过来。

伊藤佐一吓出了一身冷汗；但还没反应过来，车已经在远处停下了。他扭头一看，是那只猫挡住了那辆车。猫静悄悄地跑了过来，又过了马路，跃进了草丛。

而伊藤佐一在路边的街灯旁，看到了那枚奖章。车已经开走了；见此，他急忙冲了过去，看到街灯柔和的圆锥形灯光下，奖章早已锈迹斑斑。

伊藤佐一心痛不已；可是没等他捡起，经过的一个女孩就已经比他先捡起了那枚奖章。

\twkkai{这是……你的吗？}女孩轻声问道。

伊藤佐一这才意识到，这个女孩是同班同学，叫铃木和子\index{ling@铃木和子}；尽管她没有穿校服，他还是认出了她。

\twkkai{是的，谢谢你；}伊藤佐一接过奖章，却突然发现奖章变得焕然一新，上面刻着的东西也全都变样了。

\twkkai{这是……}

伊藤佐一低头想着，但是铃木和子已经走远。他连忙跑了上去追问。

\twkkai{你明天也来上课吗？我……我明天有事和你说；}伊藤佐一忐忑地说道。

铃木和子拨了拨她的刘海，轻声答应着。

\twkkai{可以啊——我通常都坐在最后一排的；}她补充道。

\twkkai{好的，再次谢谢你；}伊藤佐一和铃木和子告别后，便又匆匆地向家的方向走去。



